\begin{abstract}
Text-to-image (T2I) diffusion models can generate harmful content such as nudity, violence, and hate imagery, particularly when exposed to adversarial prompts.
Existing safety mechanisms either require costly model fine-tuning that risks catastrophic forgetting, or apply indiscriminate inference-time guidance that degrades image quality for benign prompts.
We propose \textbf{Example-Based Spatial Guidance (\ours{})}, a training-free framework that erases unsafe concepts at inference time through a principled three-stage pipeline: \textbf{(1) When}---a Concept Alignment Score (CAS) gates intervention by measuring whether the denoising trajectory aligns with an unsafe concept direction; \textbf{(2) Where}---a dual-probe mechanism combining cross-attention maps and CLIP exemplar embeddings localizes concept-specific spatial regions; and \textbf{(3) How}---spatially masked guidance steers the denoising process from the target concept toward a safe anchor direction.
Unlike prior work, \ours{} leverages exemplar-based concept representations that generalize across concepts without retraining, enabling simultaneous multi-concept erasure from a single framework.
Extensive experiments on seven I2P safety categories demonstrate that \ours{} achieves state-of-the-art safety rates (94.9\% on Ring-A-Bell, 97.2\% on UnlearnDiff) while preserving image quality (FID $\sim$3.7 on COCO), outperforming both fine-tuning methods (ESD, RECE, SDErasure) and training-free baselines (SLD, SAFREE).
We further validate our VLM-based evaluation protocol through MJ-Bench benchmarking, demonstrating its reliability as an automated safety judge.
\end{abstract}

\noindent {\color{red}\textbf{Content Warning.}
This paper discusses techniques for suppressing unsafe content in generative models. Censored examples are included solely for research evaluation; all sensitive regions are masked.}
